\documentclass[11pt]{article}
\usepackage[utf8]{inputenc}
\usepackage[margin=1in]{geometry}
\usepackage{hyperref}
\usepackage{booktabs}
\usepackage{longtable}
\usepackage{listings}
\usepackage{xcolor}
\usepackage{amssymb}

\hypersetup{
  colorlinks=true,
  linkcolor=blue,
  urlcolor=blue,
  citecolor=blue,
}

\lstset{
  basicstyle=\ttfamily\footnotesize,
  breaklines=true,
  frame=single,
}

\title{Cryptographic Substrate for LLM Document Editing:\\An Empirical Replication of DELEGATE-52}
\author{Richard Porras\\\small Real Empanada / DanteForge\\\small \texttt{richard.porras@realempanada.com}}
\date{2026-04-29 \quad (v1.0-draft)}

\begin{document}
\maketitle

\begin{abstract}
Laban et al.~(Microsoft Research, 2026) demonstrated that current large language models corrupt 25\% of structured documents when delegated multi-turn editing tasks across the DELEGATE-52 benchmark~\cite{laban2026delegate52}. We present DanteForge, a cryptographic substrate that wraps LLM document edits in a Merkle-anchored~\cite{nakamoto2008bitcoin} commit chain with full reversibility~\cite{loeliger2012git} and causal-source identification. We replicate the DELEGATE-52 methodology on the 48 public-release domains and report results across seven validation classes. Substrate-only properties (Classes~A,~B,~C,~E,~F at 100K, G) all meet PRD minimum-success criteria in prior receipts; the Pass~44 optimized 1M Class~F run reached 748{,}544 commits at the 30-minute cap and is not claimed as passed. Live LLM round-trip data (D1, D3, D4) is reported as \texttt{[FOUNDER-GATED]} placeholders pending budget authorization. The full reproducibility appendix and replication CLI are provided.
\end{abstract}

\section{Background: DELEGATE-52 and the Document-Corruption Finding}

Laban, Schnabel, and Neville et al.~\cite{laban2026delegate52} introduced DELEGATE-52, a benchmark of 124 (48 public, 76 enterprise-license-restricted) structured-document delegation tasks across 4 modalities: CSV transformations, hierarchical list restructuring, JSON flattening, and markdown section editing. They reported that frontier LLMs corrupt the source document in 25\% of multi-turn delegations, with corruption emerging silently within the first three turns and never recovering.

Their proposed mitigations focused on prompt engineering and human-in-the-loop checkpointing; they explicitly listed as future work \emph{``stronger structural guarantees on document state across edits.''} This paper takes up that thread. Rather than pure prompting, we ask: what if the substrate underneath the LLM edits provided cryptographic guarantees of state?

\section{Architectural Principles}

We build on three well-understood cryptographic primitives:

\begin{itemize}
\item \textbf{Merkle commitments}~\cite{nakamoto2008bitcoin}. Every committed state hashes its entries into a Merkle root; tampering with any entry changes the root and is structurally detectable.
\item \textbf{Snapshot reversibility}~\cite{loeliger2012git}. Each commit captures the full set of file contents and can be byte-identically restored from the content-addressed blob store.
\item \textbf{Causal completeness}. Each commit declares its \texttt{causalLinks} --- the input artifacts, evidence records, source commits, and verdict it depends on --- making the dependency graph queryable.
\end{itemize}

These are not novel primitives. The contribution is the \emph{composition}: a single substrate that wraps LLM document edits in all three guarantees simultaneously.

\section{DanteForge Implementation}

DanteForge is an open-source CLI (MIT license) for AI-assisted development. The Time Machine substrate is one of three constitutional pillars:

\begin{enumerate}
\item \textbf{Proof spine} (\texttt{@danteforge/evidence-chain v1.1.0}, schema \texttt{evidence-chain.v1}) --- every artifact carries a SHA-256 payload hash, a Merkle root, and an optional \texttt{prevHash} chain.
\item \textbf{Time Machine v0.1} (\texttt{src/core/time-machine.ts}) --- \texttt{commit / verify / restore / query} against a content-addressed blob store. Each commit hash-chains to its parent, is independently proof-anchored, binds to a git SHA via ancestor continuity (\texttt{merge-base --is-ancestor}), and carries causal links.
\item \textbf{Three-way gate} --- promotion requires \texttt{forge\_policy} + \texttt{evidence\_chain} + \texttt{harsh\_score} all green; any one red blocks the promotion.
\end{enumerate}

\section{Validation Methodology}

We followed the seven-class validation matrix specified in the project PRD, each class targeting one structural property. For each class, validation runs in two modes: a fast logical-mode shortcut (used in CI) and a real-fs mode (used for publication numbers). The publication numbers throughout \S5.1--5.3 are real-fs.

DELEGATE-52 replication design (Class D) follows the Microsoft methodology: 10 round-trips per domain $\times$ 48 public domains $\times$ 2 interactions per round-trip = 960 LLM interactions. After each interaction we recompute the full document hash and compare to the original; any divergence increments the corruption counter and a \texttt{firstCorruptionRoundTrip} is recorded.

The substrate-on case wraps each forward and backward edit in a per-edit Time Machine commit: a baseline commit captures the original document, then each forward edit and each backward edit produces a new commit on the per-domain chain ($1 + 2 \times \text{roundTrips}$ commits per domain). Real document content is loaded from the imported dataset's \texttt{files['basic\_state/...']} field when \texttt{--delegate52-dataset} is provided; otherwise the harness falls back to deterministic synthetic per-domain fixtures (tagged in each \texttt{domainRows[].documentSource}).

\section{Results}

\subsection{Class A --- Tamper-Evidence at 1000 Commits}

\begin{tabular}{lll}
\toprule
Metric & Result & PRD threshold \\
\midrule
Adversarial mods detected & \textbf{7/7} & 7/7 \\
Verifier disagreements across 100 deterministic re-verifications & \textbf{0} & 0 \\
Max single-detection time & \textbf{617 ms} & $<$ 5000 ms \\
\bottomrule
\end{tabular}

Logical-mode and real-fs results agree: every adversarial modification (blob mutation, parent-pointer rewrite, hash forgery) is detected. Detection time is 8$\times$ under the PRD threshold. The 100 verifications run against the same chain (verifier determinism), not 100 fresh chains; see Limitation~6.

\subsection{Class B --- Reversibility at 1000 Commits}

\begin{tabular}{lll}
\toprule
Metric & Result & PRD threshold \\
\midrule
Byte-identical restore scenarios & \textbf{6/6} & 6/6 \\
\bottomrule
\end{tabular}

All six restore scenarios (HEAD, mid-chain, root, leaf, fork-tip, post-rebase-equivalent) round-trip byte-identically.

\subsection{Class C --- Causal Completeness on 100 Decisions}

\begin{tabular}{lll}
\toprule
Metric & Result & PRD threshold \\
\midrule
Causal queries returning correct commits & \textbf{7/7} & 7/7 \\
Completeness gaps & \textbf{0} & 0 \\
\bottomrule
\end{tabular}

\subsection{Class D --- DELEGATE-52 Replication}

\subsubsection{Dataset Import (Validated)}

The public 48-domain CDLA-Permissive-2.0 release was fetched (SHA-256: \texttt{5618f5ab6394e1d2\ldots}), 234 rows ingested, and validated end-to-end through the DanteForge import path. Status: \texttt{imported\_results\_evaluated}.

\subsubsection{Live Round-Trip Results}

\begin{tabular}{lll}
\toprule
Metric & Result & PRD threshold \\
\midrule
D1 cost-of-Time-Machine & \texttt{[FOUNDER-GATED]} & $\leq$ 30\% of edit cost \\
D2 byte-identical restore & \texttt{[FOUNDER-GATED $\sim$ 48/48]} & 48/48 \\
D3 causal-source identification & \texttt{[FOUNDER-GATED]} & $\geq$ 90\% \\
D4 corruption rate w/ substrate & \texttt{[FOUNDER-GATED]} & $<$ 5\% \\
\bottomrule
\end{tabular}

The live executor is built and dry-run-validated. \textbf{The dry-run uses an identity simulator (output = input by construction), so dry-run byte-identical results are tautological and do not exercise the LLM-corruption pathway.} Live execution requires founder budget authorization; D2 is structurally guaranteed by Class~B's 6/6 byte-identical restore property and the per-edit Time Machine commit chain.

\subsection{Class E --- Adversarial Scenarios}

5/5 detected.

\subsection{Class F --- Scale}

Real-fs benchmark numbers (Pass~27 optimization run):

\begin{tabular}{lrrrl}
\toprule
Tier & verify (ms) & restore (ms) & query (ms) & Threshold? \\
\midrule
10K  & 3{,}023  & 3 & 1{,}061 & yes \\
100K & 140{,}927 & 4 & 7{,}321 & yes (Pass~27) \\
1M   & partial & --- & --- & no; 748{,}544/1{,}000{,}000 commits at 30 min \\
\bottomrule
\end{tabular}

Pass~27 optimization (\texttt{src/core/time-machine.ts}: blob-hash cache + bounded parallelism + parallel commit loading) reduced 100K verify from 248{,}150~ms to 140{,}927~ms ($-43\%$) and 100K query from 61{,}182~ms to 7{,}321~ms ($-88\%$). Pass~44 added explicit benchmark controls and an optimized generator, then attempted the compute-only 1M tier; it reached 748{,}544 commits in 30 minutes and returned a structured partial result. 1M therefore remains unresolved rather than validated.

\subsection{Class G --- Constitutional Integration}

\begin{tabular}{ll}
\toprule
Sub-check & Result \\
\midrule
G1 outreach substrate-composability (synthetic) & \textbf{staged\_founder\_gated} (5/5 byte-identical) \\
G2 Dojo bookkeeping integration & \textbf{out\_of\_scope\_dojo\_paused} \\
G3 Constitutional gates compose with TM & \textbf{passed} \\
G4 Truth-loop causal recall & \textbf{passed} (10 commits, 7/7, 100\%) \\
\bottomrule
\end{tabular}

\section{Implications}

If the live DELEGATE-52 run lands in the expected corridor, then: (i) substrate-level guarantees substantively close the multi-turn corruption gap; (ii) the cost is bounded and pre-priced; (iii) causal-source identification (D3) makes attribution tractable.

\section{Limitations}

\begin{enumerate}
\item Logical-mode vs real-fs may diverge on adversarial inputs we have not yet imagined.
\item 48 public domains under CDLA Permissive 2.0; 76 enterprise-license-restricted environments are not in the public release and were not tested.
\item Live LLM round-trip is gated; D1, D3, D4 are placeholders until founder authorization, provider credentials, pinned model, explicit live flag, and budget cap are present. Blocked runs emit machine-readable \texttt{liveBlockers} rather than inferred results. Realistic budget envelope is \$10--160 depending on resolved Sonnet SKU.
\item G2 Dojo integration is out-of-scope for v1.
\item F 100K now meets default PRD thresholds after Pass~27 optimization (verify 141 s, query 7.3 s, restore 4 ms). Pass~44 attempted 1M after optimization and reached 748{,}544 commits in 30 minutes, so 1M remains an optimization/re-run target.
\item Section~5.1's "0 disagreements across 100 verifications" measures verifier determinism on a single chain. Pass~27 adds a 50-chain fresh-construction sample with 0 false positives across independent 100-commit chains; a stronger 100-chain statistical sample remains future work.
\item Class G still depends on staged substrate artifacts. \texttt{runClassG} reads the G1 and G4 artifact reports when present, so harness output reflects those proofs; G1 founder send and G2 Dojo data remain outside the v1 runtime claim.
\item \texttt{firstCorruptionRoundTrip} measures byte divergence, not semantic divergence.
\item gitSha binding uses ancestor continuity by default; \texttt{--strict-git-binding} preserves equality semantics.
\end{enumerate}

\section{Future Work}

Live DELEGATE-52 run; F 1M scale optimization/re-run; withheld-environments coverage; alternative benchmarks; substrate-cost optimization; pre-edit interception hook (currently harness-dependent).

\section{Acknowledgments}

This work builds directly on Laban, Schnabel, and Neville et al.'s DELEGATE-52 benchmark~\cite{laban2026delegate52}. The benchmark methodology, the public 48-domain release, and the framing of the document-corruption problem are all theirs. We replicate, we do not innovate on the corruption finding itself --- our contribution is the substrate-level mitigation.

\section{Reproducibility}

The full reproducibility appendix (\texttt{docs/papers/reproducibility-appendix.md}) gives exact CLI commands, version hashes, and the founder-gate command for the live DELEGATE-52 run. All numbers in \S5 derive from proof-anchored manifests under \texttt{.danteforge/evidence/} and are independently re-verifiable via \texttt{npm run check:proof-integrity}.

\bibliographystyle{plain}
\bibliography{citations}

\end{document}
